\documentclass[a4paper]{article}     
\usepackage{amsfonts}
\usepackage{amsmath}
\usepackage{amssymb}
\usepackage{graphicx}
\usepackage{cancel}

\newcommand{\asa}{\Leftrightarrow} %als slechts als
\newcommand{\adR}{\Rightarrow} %als dan Rechts
\newcommand{\adL}{\Leftarrow}
\newcommand{\Lh}{\left(} %Links haakje
\newcommand{\Rh}{\right)}
\newcommand{\Lv}{\left[} %Links vierkanthaakje
\newcommand{\Rv}{\right]}
\newcommand{\La}{\left\{} %Links acolade
\newcommand{\Ra}{\right\}}
\newcommand{\Ras}{\right.} %Rechts acolade stelsel (omdat om stelsels te maken heb je een punt nodig
\newcommand{\pnR}{\rightarrow} %pijl naar Rechts
\newcommand{\limiet}{\lim\limits_}
\newcommand{\schrap}{\cancel}
\newcommand{\schrapb}{\bcancel} %backslash
\newcommand{\schrapk}{\xcancel} %kruis
\newcommand{\vervang}{\cancelto} 

\begin{document}

\noindent \large Auteur: Yoren Collier \\
\noindent \large Studierichting: Informatica\\
\noindent \large Oefeningenreeks 2D nr. 8.12\\

\medskip

\normalsize

$f(x) = \dfrac{\sqrt{x}-1}{\sqrt[3]{x}-1}$\\

$\limiet{x \pnR 1} f(x) = \dfrac{1-1}{1-1} = \dfrac{0}{0}$\\

$\limiet{x \pnR 1} f(x) = \limiet{x \pnR 1} 
\dfrac{\Lh \sqrt{x}-1 \Rh \Lh \sqrt{x}+1 \Rh \Lh \Lh \sqrt[3]{x} \Rh ^2 + \sqrt[3]{x} +1 \Rh} {\Lh \sqrt[3]{x}-1 \Rh \Lh \sqrt{x}+1 \Rh \Lh \Lh \sqrt[3]{x} \Rh ^2 + \sqrt[3]{x} +1 \Rh} = \limiet{x \pnR 1}
\dfrac{\schrap{\Lh x-1 \Rh} \Lh \Lh \sqrt[3]{x} \Rh ^2 + \sqrt[3]{x} +1 \Rh}{\schrap{\Lh x-1 \Rh } \Lh \sqrt{x} +1 \Rh} 
= \dfrac{1+1+1}{1+1} = \dfrac{3}{2}$\\

\begin{thebibliography}{999}
%\bibitem{a} Referentie
\end{thebibliography}
\end{document} 